\documentclass[manuscript,screen,nonacm]{acmart}

\settopmatter{printacmref=false}
\setcopyright{none}

\usepackage{pgfplotstable}
\pgfplotsset{compat=1.18}
\usepackage{booktabs}
\usetikzlibrary{patterns}

\newcommand{\fillin}[1]{\textcolor{gray}{\textit{[#1]}}}

\newread\resultstream
\newcommand{\loadresults}{
  % count data rows in the CSV (all lines minus the header), so the template
  % adapts to the number of input files
  \openin\resultstream=results/results.csv
  \def\nrows{0}
  \loop
    \read\resultstream to \resultline
    \unless\ifeof\resultstream
      \pgfmathtruncatemacro\nrows{\nrows+1}
    \repeat
  \closein\resultstream
  \pgfmathtruncatemacro\nrows{\nrows-1}
  \pgfmathtruncatemacro\lastrow{\nrows-1}
  \pgfplotstableread[col sep=comma]{results/results.csv}{\results}
  \def\tenRow{-1}\def\hundredRow{-1}\def\gigaRow{-1}\def\fourGRow{-1}
  \foreach \r in {0,...,\lastrow} {
    \pgfplotstablegetelem{\r}{filename}\of{\results}
    \ifnum\pdfstrcmp{\pgfplotsretval}{10M.txt}=0
      \ifnum\tenRow=-1 \xdef\tenRow{\r}\fi
    \fi
    \ifnum\pdfstrcmp{\pgfplotsretval}{100M.txt}=0
      \ifnum\hundredRow=-1 \xdef\hundredRow{\r}\fi
    \fi
    \ifnum\pdfstrcmp{\pgfplotsretval}{1G.txt}=0
      \ifnum\gigaRow=-1 \xdef\gigaRow{\r}\fi
    \fi
    \ifnum\pdfstrcmp{\pgfplotsretval}{4G.txt}=0
      \ifnum\fourGRow=-1 \xdef\fourGRow{\r}\fi
    \fi
  }
  \pgfplotstablegetelem{\tenRow}{ram}\of{\results}\pgfmathsetmacro\ramten{\pgfplotsretval}
  \pgfplotstablegetelem{\hundredRow}{ram}\of{\results}\pgfmathsetmacro\ramhundred{\pgfplotsretval}
  \pgfplotstablegetelem{\gigaRow}{ram}\of{\results}\pgfmathsetmacro\ramgiga{\pgfplotsretval}
  \pgfmathsetmacro\rammax{1.05*\ramgiga}
  \pgfmathsetmacro\tenAfter{\tenRow+6}
  \pgfmathsetmacro\hundredAfter{\hundredRow+6}
  \pgfmathsetmacro\gigaAfter{\gigaRow+6}
  \pgfmathsetmacro\fourGAfter{\fourGRow+6}
  % threads==8 wall time of each file (row = block start + 3 for 1,2,4,8,16,32)
  \pgfmathsetmacro\tenT{\tenRow+3}
  \pgfmathsetmacro\hundredT{\hundredRow+3}
  \pgfmathsetmacro\gigaT{\gigaRow+3}
  \pgfmathsetmacro\fourGT{\fourGRow+3}
  \pgfplotstablegetelem{\tenT}{walltime}\of{\results}\pgfmathsetmacro\wallten{\pgfplotsretval}
  \pgfplotstablegetelem{\hundredT}{walltime}\of{\results}\pgfmathsetmacro\wallhundred{\pgfplotsretval}
  \pgfplotstablegetelem{\gigaT}{walltime}\of{\results}\pgfmathsetmacro\wallgiga{\pgfplotsretval}
  \pgfplotstablegetelem{\fourGT}{walltime}\of{\results}\pgfmathsetmacro\wallfourg{\pgfplotsretval}
}

\newcommand{\loadhashes}{
  \pgfplotstableread[col sep=comma]{results/output-hash.csv}{\hashes}
}

\title{COMP90025 Project 1: Parallel Byte-Pair Encoding (BPE) Report}


\subtitle{Student Name: \fillin{Jone Doe}\\
          Student Login ID: \fillin{jdoe}
          Student ID: \fillin{12345678}}


\author{Your Name}

\begin{abstract}

This report describes a parallel OpenMP implementation of a byte-pair
encoding (BPE) training pipeline and analyses its measured performance on
\fillin{machine or cluster} across 1--32 threads for three input corpora of
10M, 100M and 1G bytes.

\fillin{You are expected to replace any content but reused the existing structure.}
\end{abstract}

\begin{document}

\maketitle

\section{Introduction}
Byte-pair encoding (BPE) is the subword-tokenization algorithm used by most
modern large language models. This project implements a parallel BPE training
pipeline in C++ with OpenMP. Task~1 reads the input file, splits it into
words, counts word frequencies, and splits each distinct word into its
characters. Task~2 repeatedly finds and merges the most frequent eligible
adjacent token pair until no eligible pair remains. The implementation is
benchmarked on \fillin{machine} with 1, 2, 4, 8, 16 and 32 OpenMP threads on
three input corpora of 10M, 100M and 1G bytes. Relative to the provided
sequential reference implementation, the parallel version reaches a speedup
of up to \fillin{speedup}$\times$ on the \fillin{largest} corpus.

\section{Results}
This section presents experimental results for the parallel BPE training pipeline. The results include wall time, peak RAM usage, and output correctness for each input corpus and number of threads.

\subsection{Measurement protocol}
All measurements were produced by the submitted \texttt{benchmark.slurm}
script on \fillin{machine, cluster, and partition}, for three input corpora
of \texttt{10M}, \texttt{100M} and \texttt{1G} bytes. For each corpus the
pipeline was run with 1, 2, 4, 8, 16 and 32 OpenMP threads,
\fillin{number of} time(s) each, and the \fillin{best or mean} wall time is
reported. Wall time is the end-to-end training time; RAM is the peak
resident set size. Output correctness is checked against the reference
output hashes in Table~\ref{tab:output-hash}.

\loadresults
\loadhashes

\subsection{Results table}
Table~\ref{tab:results} shows the measured wall time and peak RAM for each
input corpus and number of threads. The table is generated from the
\texttt{results/results.csv} file produced by the benchmark script. \fillin{Describe your observations.}
\begin{table}[h]
  \centering
  \caption{Measured BPE training results (from \texttt{results/results.csv}).}
  \label{tab:results}
  \pgfplotstabletypeset[
    col sep=comma,
    columns={filename,threads,walltime,ram},
    columns/filename/.style={string type, column type={l}, column name={File}},
    columns/threads/.style={column type={r}, column name={Threads}},
    columns/walltime/.style={column type={r}, column name={Wall time (s)}},
    columns/ram/.style={column type={r}, column name={Peak RAM (MB)}},
    every head row/.style={before row=\toprule, after row=\midrule},
    every last row/.style={after row=\bottomrule},
    font=\footnotesize,
  ]{\results}
\end{table}

\subsection{Scaling figure}
Figure~\ref{fig:scaling} shows the measured wall time versus the number of
OpenMP threads for each input corpus. The x-axis is logarithmic (base~2)
while the y-axis is linear. \fillin{Describe your observations.}

\begin{figure}[h]
  \centering
  \begin{tikzpicture}
    \begin{axis}[
        width=0.95\linewidth,
        height=0.62\linewidth,
        xlabel={Number of threads},
        ylabel={Training wall time (s)},
        grid=both, minor tick num=1,
        xmode=log, log basis x=2,
        xtick={1,2,4,8,16,32},
        legend pos=north east,
        legend entries={10M.txt, 100M.txt, 1G.txt},
      ]
      \addplot+[mark=*,thick]
        table[x=threads, y=walltime, skip coords between index={0}{\tenRow}, skip coords between index={\tenAfter}{\nrows}]{\results};
      \addplot+[mark=square*,thick]
        table[x=threads, y=walltime, skip coords between index={0}{\hundredRow}, skip coords between index={\hundredAfter}{\nrows}]{\results};
      \addplot+[mark=triangle*,thick]
        table[x=threads, y=walltime, skip coords between index={0}{\gigaRow}, skip coords between index={\gigaAfter}{\nrows}]{\results};
    \end{axis}
  \end{tikzpicture}
  \caption{Training wall time versus the number of OpenMP threads for the
  three input corpora (log2 x-axis, linear y-axis). Each line is one input
  file.}
  \label{fig:scaling}
\end{figure}

\subsection{Time versus file size (8 threads)}
Figure~\ref{fig:size} shows the wall time at a fixed thread count (8) against
the input file size, for each input corpus.
\fillin{Describe your observations.}

\begin{figure}[h]
  \centering
  \begin{tikzpicture}
    \begin{axis}[
        width=0.95\linewidth,
        height=0.5\linewidth,
        xlabel={File size (bytes)},
        ylabel={Training wall time (s)},
        xmode=log, log basis x=10,
        xtick={1e7,1e8,1e9,4e9},
        xticklabels={10M,100M,1G,4G},
        grid=both, minor tick num=1,
      ]
      % One point per input file, taken from the threads==8 row (values set
      % up by \loadresults). The x positions are the nominal file sizes.
      \addplot+[mark=*,thick]
        coordinates {(1e7,\wallten) (1e8,\wallhundred) (1e9,\wallgiga)
                     (4e9,\wallfourg)};
    \end{axis}
  \end{tikzpicture}
  \caption{Training wall time against input file size at 8 OpenMP threads
  (log10 x-axis, linear y-axis).}
  \label{fig:size}
\end{figure}

\subsection{Output hash check}
Table~\ref{tab:output-hash} shows the MD5 hash of the \texttt{output.txt}
\begin{table}[h]
  \centering
  \caption{MD5 hash of \texttt{output.txt} for each input file, compared
  against the reference output of the provided sequential baseline.}
  \label{tab:output-hash}
  \pgfplotstabletypeset[
    col sep=comma,
    columns={filename,output_md5,reference_output_md5},
    columns/filename/.style={string type, column type={l}, column name={Input file}},
    columns/output_md5/.style={string type, column type={l}, column name={Our output MD5}},
    columns/reference_output_md5/.style={string type, column type={l}, column name={Reference output MD5}},
    every head row/.style={before row=\toprule, after row=\midrule},
    every last row/.style={after row=\bottomrule},
    font=\footnotesize,
  ]{\hashes}
\end{table}

\subsection{Memory usage}
Figure~\ref{fig:ram} shows the peak RAM usage for each input corpus. \fillin{Describe your observations.}
\begin{figure}[h]
  \centering
  \begin{tikzpicture}
    \begin{axis}[
        width=0.95\linewidth,
        height=0.32\linewidth,
        ymin=0, ymax=4,
        ytick={1.5,2.5,3.5},
        yticklabels={10M.txt,100M.txt,1G.txt},
        ytick style={draw=none},
        xlabel={Peak RAM (MB)},
        xmin=0, xmax=\rammax,
        scaled x ticks=false,
        x tick label style={/pgf/number format/fixed},
      ]
      \fill[blue!30, pattern=north east lines, pattern color=blue!60]
        (axis cs:0,1.15) rectangle (axis cs:\ramten,1.85);
      \fill[red!30, pattern=north west lines, pattern color=red!60]
        (axis cs:0,2.15) rectangle (axis cs:\ramhundred,2.85);
      \fill[teal!30, pattern=crosshatch, pattern color=teal!60]
        (axis cs:0,3.15) rectangle (axis cs:\ramgiga,3.85);
    \end{axis}
  \end{tikzpicture}
  \caption{Peak RAM usage for each input corpus (measured on the 1-thread
  run). Each bar uses a different hatch pattern.}
  \label{fig:ram}
\end{figure}

\section{Analysis and Conclusion}


The results show that \fillin{summarise how training time changes with
threads, e.g. scales to a $7\times$ speedup at 8 threads on the 1G corpus}.
The speedup comes mainly from \fillin{which parallel parts of the pipeline
dominate, e.g. word counting and sorting}, while \fillin{which parts remain
sequential or bottlenecked, e.g. I/O and reductions} limits further gains.
The speedup is \fillin{below / close to} the ideal linear scaling, which is
consistent with \fillin{Amdahl's law / load imbalance / memory bandwidth}.
Peak RAM grows roughly \fillin{proportionally} with input size and is
\fillin{approximately independent of the number of threads}.

These measurements are \fillin{limits: e.g. from a single machine, a limited
range of input sizes, and few repetitions}, so they should be read as
\fillin{what they do and do not show}.

In conclusion, \fillin{one to three sentences summarising the main finding
and the most important takeaway}.

\end{document}
